\documentclass[aspectratio=169]{beamer}
\input{header}
\coursecode{JKSSB}

\title{Bits, Bytes \& Units}
\subtitle{Lesson 03 --- Measuring Digital Information}
\author{JKSSB Computer Awareness}
\date{\today}

\begin{document}
\frame{\titlepage}

\begin{frame}{Video Lesson 03}
  \centering
  \Large\href{https://youtu.be/j_62CGCBElo}{Watch: Bits, Bytes \& KB--TB}
\end{frame}

\begin{frame}{Why Units Matter}
  \begin{itemize}
    \item A file is reported as 4 MB; a network plan is reported as 20 Mbps.
    \item Both describe digital data, but one uses bytes and the other uses
      bits per second.
    \item Mixing the unit or the case of the letter can create an eightfold
      mistake.
  \end{itemize}
\end{frame}

\begin{frame}{The Small Building Blocks}
  \begin{itemize}
    \item \key{Bit}: a binary digit, 0 or 1; the smallest unit of data.
    \item \key{Nibble}: 4 bits.
    \item \key{Byte}: 8 bits (two nibbles).
    \item \key{Word}: the natural data size a processor handles as a unit;
      its size depends on the architecture.
    \item \key{Character}: a text symbol; its encoded storage is not always
      one byte.
  \end{itemize}
\end{frame}

\begin{frame}{Case Is Part of the Unit}
  \begin{block}{Do not confuse}
    \texttt{b} means bit; \texttt{B} means byte. One byte contains 8 bits.
  \end{block}
  \begin{itemize}
    \item \texttt{Mbps}: megabits per second.
    \item \texttt{MB}: megabytes.
    \item At matching decimal prefixes, 8 Mbps corresponds to at most
      1 MB/s before protocol overhead.
  \end{itemize}
\end{frame}

\begin{frame}{Capacity Ladder: Exam Convention}
  \begin{center}
    \begin{tabular}{ll}
      \toprule
      \textbf{Conversion} & \textbf{JKSSB-style binary convention} \\
      \midrule
      1 KB & 1024 bytes \\
      1 MB & 1024 KB \\
      1 GB & 1024 MB \\
      1 TB & 1024 GB \\
      \bottomrule
    \end{tabular}
  \end{center}
  \vspace{0.25cm}
  \muted{Use 1024 when an exam question uses the traditional KB/MB/GB/TB
  ladder. Modern standards also distinguish decimal SI units from binary
  IEC units such as KiB.}
\end{frame}

\begin{frame}{One-Step Conversion}
  \begin{exampleblock}{Example}
    How many bytes are in 3 KB under the 1024 convention?
  \end{exampleblock}
  \begin{enumerate}
    \item Write the given unit: 3 KB.
    \item Use 1 KB = 1024 bytes.
    \item Calculate: $3 \times 1024 = 3072$ bytes.
  \end{enumerate}
  \begin{alertblock}{Check}
    State the convention and unit with the answer.
  \end{alertblock}
\end{frame}

\begin{frame}{Capacity Is Not Speed}
  \begin{itemize}
    \item \key{Capacity} tells how much data can be stored.
    \item \key{Transfer rate} tells how much data moves in a time interval.
    \item Storage is commonly labelled in bytes; data rates are commonly
      labelled in bits per second.
    \item A large-capacity drive is not automatically a fast drive.
  \end{itemize}
\end{frame}

\begin{frame}{Lesson 03: Recall}
  \begin{itemize}
    \item 1 nibble = 4 bits; 1 byte = 8 bits.
    \item The bit/byte distinction depends on lowercase \texttt{b} vs
      uppercase \texttt{B}.
    \item A processor's word size is architecture-dependent.
    \item Under the exam's 1024 convention: 1 KB = 1024 bytes.
  \end{itemize}
\end{frame}
\end{document}
